% ============================================================
%  Números Racionais — Apostila Premium | PratiqueJá
%  Engine: XeLaTeX
% ============================================================
\documentclass[12pt]{article}
\usepackage{../../pratiqueja}
\usepackage{cancel}

% \hlm — destaque de resultado em modo-math (cor sem bold)
\newcommand{\hlm}[1]{{\color{darkblue}#1}}

\setsubject{Números Racionais}

\begin{document}
\pagenumbering{arabic}
\setstretch{1.2}
\color{bodytext}

\heroheader{Números Racionais}%
           {Frações, Operações e Dízimas Periódicas}%
           {Matemática · Básico}

\setlength{\columnsep}{18pt}
\setlength{\columnseprule}{0.5pt}
\def\columnseprulecolor{\color{exborder}}

\begin{multicols}{2}

%─────────────────────────────────────────────────────────────
\TW{Def}{badgepurple}{Conjunto dos Racionais}{Frações, decimais e inteiros}

\formula{$\displaystyle\mathbb{Q} = \left\{ \frac{a}{b} \;\middle|\; a \in \mathbb{Z},\ b \in \mathbb{Z}^{*} \right\}$}

\small\color{bodytext}%
Todo número racional é escrito como $\frac{a}{b}$: $a$ é o \textbf{numerador}
e $b$ o \textbf{denominador} ($b\neq0$). Ao dividir, o resultado pode ser:

\exemp{%
  \textbf{Inteiro:}\quad
    $\dfrac{6}{2}=3$\qquad $\dfrac{-10}{5}=-2$\\[3pt]
  \textbf{Decimal exato:}\quad
    $\dfrac{1}{2}=0{,}5$\qquad $\dfrac{3}{4}=0{,}75$\\[3pt]
  \textbf{Dízima simples} (período logo após a vírgula):\\
  \quad $\dfrac{1}{3}=0{,}333\ldots=0{,}\overline{3}$\quad
        $\dfrac{7}{11}=0{,}\overline{63}$\\[3pt]
  \textbf{Dízima composta} (parte não periódica antes):\\
  \quad $\dfrac{7}{12}=0{,}58\overline{3}$\quad
        $\dfrac{5}{6}=0{,}8\overline{3}$%
}

%─────────────────────────────────────────────────────────────
\TW{Tipos}{badgegreen}{Tipos de Fração}{Própria, imprópria, aparente e mista}

\exemp{%
  \textbf{Própria} — num $<$ den, resultado entre 0 e 1:\\
  \quad $\dfrac{2}{6},\quad\dfrac{3}{4},\quad\dfrac{5}{9}$\\[3pt]
  \textbf{Imprópria} — num $>$ den, resultado $>1$:\\
  \quad $\dfrac{7}{2},\quad\dfrac{5}{3},\quad\dfrac{8}{5}$\\[3pt]
  \textbf{Aparente} — numerador múltiplo do denominador:\\
  \quad $\dfrac{8}{2}=4,\quad\dfrac{6}{3}=2,\quad\dfrac{12}{4}=3$\\[3pt]
  \textbf{Mista} — parte inteira + fração própria:\\
  \quad $2\dfrac{3}{5}=\dfrac{2\times5+3}{5}=\hlm{\dfrac{13}{5}}$%
}

%─────────────────────────────────────────────────────────────
\TW{=}{darkblue}{Frações Equivalentes}{Amplificação e simplificação}

\regra{$\dfrac{a}{b} = \dfrac{a\times k}{b\times k}$ \quad $(k\neq0)$ —
multiplicar ou dividir num. e den. por $k$ não altera o valor.}

\SH{Amplificação — multiplicar por $k$}
\exemp{$\dfrac{1}{3}=\dfrac{2}{6}=\dfrac{4}{12}$ — todas equivalentes}

\SH{Simplificação — dividir pelo MDC}
\exemp{%
  $\dfrac{12}{18}$: MDC$(12,18)=6$
    $\;\Rightarrow\;\dfrac{12\div6}{18\div6}=\hlm{\dfrac{2}{3}}$\\[3pt]
  $\dfrac{30}{45}$: MDC$(30,45)=15$
    $\;\Rightarrow\;\dfrac{30\div15}{45\div15}=\hlm{\dfrac{2}{3}}$%
}
\obs{Fração \textbf{irredutível}: MDC(numerador, denominador) = 1.}

%─────────────────────────────────────────────────────────────
\TW{Conv}{badgegreen}{Conversão: Mista $\leftrightarrow$ Imprópria}{Divisão euclidiana}

\SH{Mista $\to$ Imprópria}
\formula{$n\,\dfrac{p}{q} = \dfrac{n\times q + p}{q}$}

\exemp{%
  $3\dfrac{2}{5}=\dfrac{3\times5+2}{5}=\hlm{\dfrac{17}{5}}$\\[3pt]
  $4\dfrac{1}{3}=\dfrac{4\times3+1}{3}=\hlm{\dfrac{13}{3}}=4{,}333\ldots$%
}

\SH{Imprópria $\to$ Mista (divisão euclidiana)}
\exemp{%
  $\dfrac{17}{5}$: $17\div5=3$ (resto $2$)
    $\;\Rightarrow\;\hlm{3\dfrac{2}{5}}$\\[3pt]
  $\dfrac{22}{7}$: $22\div7=3$ (resto $1$)
    $\;\Rightarrow\;\hlm{3\dfrac{1}{7}}$%
}

%─────────────────────────────────────────────────────────────
\TW{$<\!>$}{darkblue}{Comparação de Frações}{Qual fração é maior?}

\SH{Mesmo denominador — compare os numeradores}
\exemp{%
  $\dfrac{3}{7}$ vs $\dfrac{5}{7}$: $3<5$
    $\;\Rightarrow\;\hlm{\dfrac{3}{7}<\dfrac{5}{7}}$%
}

\SH{Produto cruzado — método rápido}
\regra{$\dfrac{a}{b}\;\square\;\dfrac{c}{d}
       \;\Leftrightarrow\; a\times d \;\square\; c\times b$}

\exemp{%
  $\dfrac{3}{4}$ vs $\dfrac{5}{6}$:
    $3\times6=18$,\; $5\times4=20$\\
  $18<20\;\Rightarrow\;\hlm{\dfrac{3}{4}<\dfrac{5}{6}}$\\[3pt]
  $\dfrac{7}{9}$ vs $\dfrac{5}{6}$:
    $7\times6=42$,\; $5\times9=45$\\
  $42<45\;\Rightarrow\;\hlm{\dfrac{7}{9}<\dfrac{5}{6}}$%
}

\SH{Denominadores diferentes — MMC}
\exemp{%
  $\dfrac{3}{4}$ vs $\dfrac{5}{6}$ — MMC$(4,6)=12$:\\
  $\hlm{\dfrac{9}{12}}$ vs $\hlm{\dfrac{10}{12}}$
    $\;\Rightarrow\;\dfrac{3}{4}<\dfrac{5}{6}$%
}

%─────────────────────────────────────────────────────────────
\TW{D\'iz}{badgepurple}{Dízima Periódica $\to$ Fração}{Técnica da equação: eliminar o período}

\small\color{bodytext}%
Chame a dízima de $x$. Multiplique por potência de 10 que alinhe os
períodos e subtraia.

\SH{Dízima simples — período logo após a vírgula}
\exemp{%
  $x=0{,}\overline{3}$\;→\; $10x=3{,}333\ldots$\\
  $10x-x=3$\;→\; $9x=3$
    \;→\; $\hlm{x=\dfrac{3}{9}=\dfrac{1}{3}}$\\[4pt]
  $x=0{,}\overline{63}$\;→\; $100x=63{,}6363\ldots$\\
  $100x-x=63$\;→\; $99x=63$
    \;→\; $\hlm{x=\dfrac{63}{99}=\dfrac{7}{11}}$%
}

\SH{Dízima composta — parte não periódica antes}
\exemp{%
  $x=0{,}58\overline{3}$ (1 não-per., 1 per.)\\
  $10x=5{,}8\overline{3}$\quad $100x=58{,}\overline{3}$\\
  $90x=52{,}5\;\Rightarrow\;900x=525$
    \;→\; $\hlm{x=\dfrac{7}{12}}$%
}
\nota{\textbf{Atalho:}\;
$\dfrac{\text{com período}-\text{sem período}}%
       {\underbrace{9\cdots9}_{\text{peri.}}%
        \underbrace{0\cdots0}_{\text{não peri.}}}$\quad
Ex.: $0{,}58\overline{3}=\dfrac{583-58}{900}=\dfrac{7}{12}$}

%─────────────────────────────────────────────────────────────
\TW{$+$/$-$}{darkblue}{Soma e Subtração de Frações}{Denominadores iguais e diferentes}

\SH{Mesmo denominador — some/subtraia os numeradores}
\formula{$\dfrac{a}{c}\pm\dfrac{b}{c}=\dfrac{a\pm b}{c}$}
\exemp{%
  $\dfrac{2}{4}+\dfrac{1}{4}=\hlm{\dfrac{3}{4}}$\qquad
  $\dfrac{2}{4}-\dfrac{1}{4}=\hlm{\dfrac{1}{4}}$%
}

\SH{Método MMC — denominadores diferentes}
\exemp{%
  $\dfrac{2}{3}+\dfrac{1}{4}$ — MMC$(3,4)=12$\\
  $\dfrac{8}{12}+\dfrac{3}{12}=\hlm{\dfrac{11}{12}}$%
}

\SH{Método cruzado — mais rápido para 2 frações}
\regra{$\dfrac{a}{b}\pm\dfrac{c}{d}=\dfrac{(a\times d)\pm(c\times b)}{b\times d}$}
\exemp{%
  \textbf{Soma:}\;
    $\dfrac{2}{3}+\dfrac{1}{4}=\dfrac{8+3}{12}=\hlm{\dfrac{11}{12}}$\\[3pt]
  \textbf{Subtração:}\;
    $\dfrac{5}{2}-\dfrac{3}{7}=\dfrac{35-6}{14}=\hlm{\dfrac{29}{14}}$%
}
\nota{\textbf{Quando usar cada método:} o MMC é mais eficiente quando
os denominadores têm fatores comuns (ex.: 4 e 6 → MMC=12, não 24).
O cruzado é direto quando são primos entre si.}

%─────────────────────────────────────────────────────────────
\TW{$\times$/$\div$}{darkblue}{Multiplicação e Divisão}{Num.$\times$num., den.$\times$den.; inverso na divisão}

\SH{Multiplicação — direto}
\formula{$\dfrac{a}{b}\times\dfrac{c}{d}=\dfrac{a\times c}{b\times d}$}
\exemp{%
  $\dfrac{2}{3}\times\dfrac{3}{4}=\dfrac{6}{12}=\hlm{\dfrac{1}{2}}$\qquad
  $\dfrac{5}{2}\times\dfrac{3}{7}=\hlm{\dfrac{15}{14}}$%
}

\SH{Divisão — multiplica pelo inverso}
\formula{$\dfrac{a}{b}\div\dfrac{c}{d}=\dfrac{a}{b}\times\dfrac{d}{c}=\dfrac{a\times d}{b\times c}$}
\exemp{%
  $\dfrac{2}{3}\div\dfrac{4}{6}
    =\dfrac{2}{3}\times\dfrac{6}{4}=\dfrac{12}{12}=\hlm{1}$\\[3pt]
  $\dfrac{3}{4}\div\dfrac{8}{5}
    =\dfrac{3}{4}\times\dfrac{5}{8}=\hlm{\dfrac{15}{32}}$%
}
\nota{\textbf{Simplifique antes de multiplicar} quando possível —
cancele fatores entre numeradores e denominadores de frações
\emph{diferentes} antes de efetuar a conta:\\
$\dfrac{2}{\cancel{3}}\times\dfrac{\cancel{3}}{4}
 =\dfrac{2}{1}\times\dfrac{1}{4}=\hlm{\dfrac{1}{2}}$}

%─────────────────────────────────────────────────────────────
\TW{$a^n$}{darkblue}{Potenciação de Frações}{Eleva numerador e denominador ao expoente}

\formula{$\left(\dfrac{a}{b}\right)^{\!n}=\dfrac{a^n}{b^n}$
         \qquad
         $\left(\dfrac{a}{b}\right)^{\!-n}=\dfrac{b^n}{a^n}$}

\exemp{%
  $\left(\dfrac{2}{3}\right)^3=\hlm{\dfrac{8}{27}}$\qquad
  $\left(\dfrac{3}{4}\right)^2=\hlm{\dfrac{9}{16}}$\\[3pt]
  \textbf{Expoente zero:}\;
    $\left(\dfrac{5}{7}\right)^0=\hlm{1}$
    — qualquer fração não nula\\[3pt]
  \textbf{Expoente negativo:}\;
    $\left(\dfrac{2}{3}\right)^{-2}
    =\left(\dfrac{3}{2}\right)^2=\hlm{\dfrac{9}{4}}$%
}

%─────────────────────────────────────────────────────────────
\TW{Mista}{badgegreen}{Operações com Frações Mistas}{Converta para imprópria antes de operar}

\small\color{bodytext}%
O procedimento padrão é \textbf{converter as mistas em impróprias},
operar e, se necessário, converter o resultado de volta.

\exemp{%
  \textbf{Soma:}\;
    $2\dfrac{1}{3}+1\dfrac{3}{4}=\dfrac{7}{3}+\dfrac{7}{4}$\\
  MMC$(3,4)=12$:\;
    $\dfrac{28}{12}+\dfrac{21}{12}=\dfrac{49}{12}=\hlm{4\dfrac{1}{12}}$\\[4pt]
  \textbf{Multiplicação:}\;
    $1\dfrac{1}{2}\times2\dfrac{2}{3}=\dfrac{3}{2}\times\dfrac{8}{3}$\\
  $=\dfrac{\cancel{3}}{2}\times\dfrac{8}{\cancel{3}}=\dfrac{8}{2}=\hlm{4}$%
}
\nota{Em somas de mistas com denominadores iguais é possível somar
as partes inteiras separadamente e as fracionárias separadamente,
reagrupando ao final. Mesmo assim, a conversão para imprópria
é sempre a abordagem mais segura.}

%─────────────────────────────────────────────────────────────
\TW{Prob}{badgegreen}{Problemas Contextualizados}{Receitas, áreas, velocidade e descontos}

\exemp{%
  \textbf{Receita:} pede $\dfrac{3}{4}$ xíc.; fazer $\dfrac{1}{2}$ da receita:\\
  $\dfrac{3}{4}\times\dfrac{1}{2}=\hlm{\dfrac{3}{8}}$ xícara\\[4pt]
  \textbf{Área:} jardim $2\dfrac{1}{2}\,$m $\times$ $1\dfrac{3}{4}\,$m:\\
  $\dfrac{5}{2}\times\dfrac{7}{4}=\dfrac{35}{8}=\hlm{4\dfrac{3}{8}\,\text{m}^2}$\\[4pt]
  \textbf{Velocidade:} $\dfrac{5}{2}$\,km em $\dfrac{1}{4}$\,h:\\
  $v=\dfrac{5}{2}\div\dfrac{1}{4}=\dfrac{5}{2}\times4=\hlm{10\,\text{km/h}}$\\[4pt]
  \textbf{Desconto:} R\$\,120 com $\dfrac{1}{5}$ de desconto:\\
  $120\times\dfrac{1}{5}=24\;\Rightarrow$ preço final: \hl{R\$\,96}%
}

%─────────────────────────────────────────────────────────────
\TW{!}{vered}{Erros Comuns}{Armadilhas frequentes com frações}

\exemp{%
  \textbf{Somar denominadores \nok}\\
  $\dfrac{1}{2}+\dfrac{1}{3}\neq\dfrac{2}{5}$\quad
  Correto: $\dfrac{3+2}{6}=\hlm{\dfrac{5}{6}}$\\[4pt]
  \textbf{Não inverter na divisão \nok}\\
  $\dfrac{2}{3}\div\dfrac{4}{5}\neq\dfrac{8}{15}$\quad
  Correto: $\dfrac{2}{3}\times\dfrac{5}{4}=\hlm{\dfrac{5}{6}}$\\[4pt]
  \textbf{Deixar sem simplificar \nok}\\
  $\dfrac{6}{9}$ incompleto $\;\Rightarrow\;$ MDC$=3$
    $\;\Rightarrow\;\hlm{\dfrac{2}{3}}$\\[4pt]
  \textbf{Operar fração mista diretamente \nok}\\
  Converta sempre:
  $2\dfrac{1}{3}\times3=\dfrac{7}{3}\times3=\hlm{7}$%
}

%─────────────────────────────────────────────────────────────
\TW{$\mathbb{R}$}{badgegreen}{Irracionais e Reais}{Dízimas não periódicas e o conjunto $\mathbb{R}$}

\regra{Um número \textbf{irracional} \emph{não} pode ser escrito como
$\dfrac{a}{b}$ com $a,b\in\mathbb{Z}$ e $b\neq0$. Em decimal:
dízima \textbf{não periódica e infinita}.}

\SH{Irracionais famosos}
\exemp{%
  $\sqrt{2}=1{,}41421\ldots$\quad $\pi=3{,}14159\ldots$\\
  $e=2{,}71828\ldots$\quad
  $\sqrt{3},\;\sqrt{5},\;\sqrt{7},\ldots$%
}

\SH{Diagrama de inclusão}
\formula{$\mathbb{N}\subset\mathbb{Z}\subset\mathbb{Q}\subset\mathbb{R}$}

\exemp{%
  $\sqrt{9}=3$\;→\; \hl{natural, inteiro, racional, real}\\
  $\sqrt{2}$\;→\; \hl{irracional (real, mas não racional)}\\
  $-\dfrac{5}{3}$\;→\; \hl{racional e real, não inteiro}%
}
\nota{\textbf{Como identificar irracionais:} $\sqrt{n}$ é irracional
sempre que $n$ não for um quadrado perfeito.\\
Quadrados de 1 a 100: 1, 4, 9, 16, 25, 36, 49, 64, 81, 100.}

\end{multicols}
\end{document}
